\documentclass[11pt]{scrartcl}
%%
%
% This is a poster template with latex macros and using
% the University of Florida Logo.  For further information
% on making postscript, resizeing, and printing the poster file
% please see web site 
% http://www.phys.ufl.edu/~pjh/posters/poster_howto.html
% 
% N.B. This format is cribbed from one obtained from the University
% of Karlsruhe, so some macro names and parameters are in German
% Here is a short glosary:
% Breite: width
% Hoehe: height
% sPALTE: column
% Kasten: box
%
% All style files necessary are part of standard TeTeX distribution
% On the UF unix cluster you should not need to import these files
% specially, as they will be automatically located.  If you
% run on a local PC however, you will need to locate these files.
% At UF try /usr/local/TeTeX...
% 
% P. Hirschfeld 2/11/00
%
% The recommended procedure is to first generate a ``Special Format" size poster
% file, which is relatively easy to manipulate and view.  It can be
% resized later to A0 (900 x 1100 mm) full poster size, or A4 or Letter size
% as desired (see web site).  Note the large format printers currently
% in use at UF's OIR have max width of about 90cm or 3 ft., but the paper
% comes in rolls so the length is variable.  See below the specifications
% for width and height of various formats.  Default in the template is
% ``Special Format",  with 4 columns.
%%
%% 
%% Choose your poster size:
%% For printing you will later RESIZE your poster by a factor
%%        2*sqrt(2) = 2.828    (for A0)
%%        2         = 2.00     (for A1) 
%%  
%% 
\def\breite{400mm}     % Special Format. 
\def\hoehe{400mm}      % Scaled by 2.82 this gives 110cm x 90cm 

%\def\breite{390mm}     % Special Format. 
%\def\hoehe{319.2mm}      % Scaled by 2.82 this gives 110cm x 90cm 
\def\anzspalten{2}
%%
%%\def\breite{420mm}     % A3 LANDSCAPE
%%\def\hoehe{297mm}
%%\def\anzspalten{4}
%%
%% \def\breite{297mm}     % A3 PORTRAIT
%% \def\hoehe{420mm}
%% \def\anzspalten{3}
%%
%% \def\breite{210mm}     % A4 PORTRAIT
%% \def\hoehe{297mm}
%% \def\anzspalten{2}
%%
%%
%% Procedure:
%%   Generate poster.dvi with latex
%%   Check with Ghostview
%%   Make a .ps-file with ``dvips -o poster.ps poster''
%%   Scale it with poster_resize poster.ps S
%%   where S is scale factor
%%     for Special Format->A0 S= 2.828 (= 2^(3/2)))
%%     for Special Format->A1 S= 2 (= 2^(2/2)))
%% 
%% Sizes (European:)
%%   A3: 29.73 X 42.04 cm
%%   A1: 59.5 X 84.1 cm
%%   A0: 84.1 X 118.9 cm
%%   N.B. The recommended procedure is ``Special Format x 2.82"
%%   which gives 90cm x 110cm (not quite A0 dimensions).
%%
%% --------------------------------------------------------------------------
%%
%% Load the necessary packages
%% 
\usepackage{palatino}
\usepackage[latin1]{inputenc}
\usepackage{epsf}
\usepackage{graphicx,psfrag,color,pstricks,pst-grad}
\usepackage{amsmath,amssymb}
\usepackage{latexsym}
\usepackage{calc}
\usepackage{multicol}
\usepackage{url}
%%
%% Define the required numbers, lengths and boxes 
%%
\newsavebox{\dummybox}
\newsavebox{\spalten}
%\input psfig.sty

%%
%%
\newlength{\bgwidth}\newlength{\bgheight}
\setlength\bgheight{\hoehe} \addtolength\bgheight{-1mm}
\setlength\bgwidth{\breite} \addtolength\bgwidth{-1mm}

\newlength{\kastenwidth}

%% Set paper format
\setlength\paperheight{\hoehe}                                             
\setlength\paperwidth{\breite}
\special{papersize=\breite,\hoehe}

\topmargin -1in
\marginparsep0mm
\marginparwidth0mm
\headheight0mm
\headsep0mm


%% Minimal Margins to Make Correct Bounding Box
%\setlength{\oddsidemargin}{-2.44cm}
%\addtolength{\topmargin}{-3mm}
\setlength{\oddsidemargin}{-2.44cm}
\addtolength{\topmargin}{-3mm}
\textwidth\paperwidth
\textheight\paperheight

%%
%%
\parindent0cm
\parskip1.5ex plus0.5ex minus 0.5ex
\pagestyle{empty}




\definecolor{recoilcolor}{rgb}{1,0,0}
\definecolor{occolor}{rgb}{0,1,0}
\definecolor{pink}{rgb}{0,1,1}





\def\UberStil{\normalfont\sffamily\bfseries\large}
\def\UnterStil{\normalfont\sffamily\small}
\def\LabelStil{\normalfont\sffamily\tiny}
\def\LegStil{\normalfont\sffamily\tiny}

%%
%% Define some commands
%%
\definecolor{JG}{rgb}{0.1,0.9,0.3}

\newenvironment{kasten}{
  \begin{lrbox}{\dummybox}
    \begin{minipage}{\linewidth}}
    {\end{minipage}
  \end{lrbox}
  \raisebox{-\depth}{\psshadowbox[cornersize=absolute,linearc=14pt,framesep=1em]{\usebox{\dummybox}}}\\[0.5em]}
\newenvironment{spalte}{
  \setlength\kastenwidth{1.2\textwidth}
  \divide\kastenwidth by \anzspalten
  \begin{minipage}[t]{\kastenwidth}}{\end{minipage}}

%\renewcommand{\emph}[1]{{\color{red}\textbf{#1}}}


%\def\op#1{\hat{#1}}
\begin{document}
%%%%%%%%%%%%%%%%%%%%%%%%%%%%%%%%%%%%%%%%%%%%%%%%%%%%
%%%               Background                     %%%
%%%%%%%%%%%%%%%%%%%%%%%%%%%%%%%%%%%%%%%%%%%%%%%%%%%%
{\newrgbcolor{gradbegin}{0.5 0.5 1}
  \newrgbcolor{gradend}{1 1 1}%{1 1 0.5}%
  %\rput[cm](15.23,-15){
  \rput[tl]{0}(0,0){
  \includegraphics[width=\paperwidth]{logos/antena1} %lower resolution
  %\includegraphics[width=\paperwidth]{logos/alma_logo} %higher resolution
  }
\vfill}

%{\newrgbcolor{gradbegin}{0.5 0.5 1}%
%  \newrgbcolor{gradend}{1 1 1}%{1 1 0.5}%
%  \psframe[fillstyle=gradient,gradend=gradend,%
%  gradbegin=gradbegin,gradmidpoint=0.1](\bgwidth,-\bgheight)}
%\vfill
%%%%%%%%%%%%%%%%%%%%%%%%%%%%%%%%%%%%%%%%%%%%%%%%%%%%
%%%                     Header                   %%%
%%%%%%%%%%%%%%%%%%%%%%%%%%%%%%%%%%%%%%%%%%%%%%%%%%%%
\hfill
\psshadowbox{\makebox[0.95\textwidth]{%
    \hfill
	\parbox[c]{2cm}{\includegraphics[width=2cm,height=!]{logos/alma_logo}}
    \hfill
    \parbox[c]{0.7\linewidth}{%
      \begin{center}
		\textbf{\Huge {Title}}\\[0.5em]
		\textsc{\large FirstName LastName$^{1}$}
		\\[0.3em]
		{ $^1$Universidad T\'ecnicaFederico Santa Mar\'ia, Valpara\'iso, Chile}
      \end{center}}
	\hfill
    \parbox[c]{2cm}{\includegraphics[width=3cm,height=!]{logos/usm_logo}}
	\hfill
\hfill}}\hfill\mbox{}\\

\hfill
\psshadowbox{\makebox[0.83\textwidth]{%
    \hfill
    \parbox[t]{0.8\linewidth}{%
\hfill{\large\bf{\color{red}ABSTRACT}}\hfill\mbox{}\\
}
\hfill}}\hfill\mbox{}

\begin{lrbox}{\spalten}
  \parbox[t][0.6\textheight]{1.3\textwidth}{
   % \vspace*{0.2cm}
%%%%%%%%%%%%%%%%%%%%%%%%%%%%%%%%%%%%%%%%%%%%%%%%%%%%
%%%                 first column                 %%%             
%%%%%%%%%%%%%%%%%%%%%%%%%%%%%%%%%%%%%%%%%%%%%%%%%%%%
    \begin{spalte}
     
	  \begin{kasten}
        \section*{\hspace{0.2cm} {\color{red} Topic 1} }
			\begin{minipage}[t]{0.95\linewidth}
		\end{minipage} \hfill
      \end{kasten}
	  \begin{kasten}
        \section*{\hspace{0.2cm} {\color{red} Topic 2} }
			\vspace{-0.1cm}
			\begin{minipage}[t]{0.43\linewidth}
			\end{minipage}
	  \end{kasten}
    \end{spalte}
	 \hfill
%%%%%%%%%%%%%%%%%%%%%%%%%%%%%%%%%%%%%%%%%%%%%%%%%%%%
%%%               second column                  %%%             
%%%%%%%%%%%%%%%%%%%%%%%%%%%%%%%%%%%%%%%%%%%%%%%%%%%%
    \begin{spalte}
       
      \begin{kasten}
        \section*{\hspace{0.2cm} {\color{red} Topic 3} }

      \end{kasten}

    \end{spalte}
     \hfill\mbox{}
}

\end{lrbox}
\hfill
\resizebox*{0.95\textwidth}{!}{
  \usebox{\spalten}}\hfill\mbox{}

\vspace{0.1cm}
\hspace{1cm}
\psshadowbox[cornersize=absolute,linearc=14pt]{\makebox[0.923\textwidth]{%
 \hfill
 \parbox[t]{0.9\linewidth}{
\begin{minipage}[t]{0.95\linewidth}
\section*{ {\normalsize \color{red} Conclusions}}
{\footnotesize
%\renewcommand{\baselinestretch}{0.5}
			\begin{minipage}[t]{0.50\linewidth}
			\end{minipage} \hspace{0.2cm}

			\begin{minipage}[t]{0.50\linewidth}
			\end{minipage}

}
\end{minipage}
}\hfill
}}\hfill\mbox{}\\
\vfill

\hfill
\psshadowbox{\makebox[0.95\linewidth]{
\hfill
\begin{minipage}[t]{0.30\linewidth}
	{\scriptsize
	{\small\bf Acknowledgements}
	}
	\renewcommand{\baselinestretch}{0.5}
\end{minipage}\hfill

\begin{minipage}[t]{0.60\linewidth}
	{\scriptsize
	\renewcommand{\baselinestretch}{1.0}
	{\small\bf References}
        [1] Graham, Matthew and Harrison, Paul and Morris, Dave and Rixon, ``VOSpace service specification, Version 1.01'' (2007).
        \renewcommand{\baselinestretch}{0.5}
	}
\end{minipage}
}}\hfill\mbox{}

\hfill
\psshadowbox{\makebox[0.95\linewidth]{
\begin{minipage}[t]{0.15\linewidth}
	\begin{tabular}{cc}
	\includegraphics[width=!,height=1.5cm]{logos/alma} &
	\includegraphics[width=!,height=1.5cm]{logos/utfsm}
	\end{tabular}
\end{minipage} \hfill
\begin{minipage}[t]{0.50\linewidth}
	\begin{center}
	\Huge{Computer System Research Group}
	\end{center}
\end{minipage} \hfill
\begin{minipage}[t]{0.18\linewidth}
	\begin{tabular}{cccc}
	\end{tabular}
\end{minipage}
}}\hfill\mbox{}


\vfill
\hfill
{\scriptsize
Contact e-mail: {\color{blue}user@csrg.cl} / For more information about
Computer System Researc Group, please visit our web site
{\color{blue}http://www.csrg.cl}
}
\hfill
\

\end{document}
